% RESEARCH_DRAFT_SHA256: d58d649cb31ec17138eb763a23a9894b770ecc4e46d8321a27fe144986c24518
% Layout and source-reference wrappers expanded for the web renderer.
\documentclass[11pt]{article}
\usepackage[margin=25mm]{geometry}
\usepackage{fontspec,unicode-math}
\setmainfont{Latin Modern Roman}
\setmathfont{Latin Modern Math}
\usepackage{amsmath,amssymb,mathtools,xurl,hyperref}
\hypersetup{colorlinks=true,urlcolor=blue,linkcolor=blue}
\setlength{\parindent}{0pt}
\setlength{\parskip}{5pt}
\setlength{\emergencystretch}{3em}
\newcommand{\sref}[2]{\hyperlink{source:#1:#2}{[S#2]}}
\newcommand{\cataloguescope}[1]{\paragraph{Source question.}#1}
\newcommand{\floor}[1]{\left\lfloor#1\right\rfloor}
\begin{document}

\section*{Dimension witnesses: amplification invariance and exact anticommutation relations}
\noindent catalogue entry; catalogue entry. Research attempt dated 11 October 2026.
\paragraph{Result in scope.}This manuscript contains a source-specific attempt and complete proofs of the stated auxiliary calculations. It makes no new solution or novelty claim for the entire target.
\section{Original problem and preserved definition}
The following definition and source list are reproduced from the repository. Any earlier inspection dates and reported literature claims in this preserved text describe that source record, not additional verification by this attempt.
% BEGIN PRESERVED DEFINITION
\subsection{Definitions and mathematical statement}

\paragraph{Shared notation.}Standard number systems are $\mathbb N,\mathbb Z,\mathbb Q,\mathbb R,\mathbb C$. Specialized notation and index ranges follow the cited source and the definitions below.


A quantum function is represented by the source operator-valued measurements on a Hilbert space. Test query complexity counts calls to that representation; robustness measures how test performance changes with deviations from the required property.
\paragraph{Record scope and qualifications.}
The catalogue title is retained as metadata. Named certificates, contracts, bounds or frameworks in that title are not asserted as results of the original source extract.
The supplied registry flags this extraction as unrecoverable. It is displayed as a literal source fragment; the definitions provide context rather than a replacement formal statement.

It also includes an incomplete stated theorem whose final condition is missing. The fragment is shown literally rather than supplied with an invented conclusion.

Source-record limitation: Source flags the statement as unrecoverable: The linked report was inspected: the apparent question is immediately answered by the stated corollary of $\mathrm{MIP}^*=\mathrm{RE}$, so the extraction is not an unresolved problem.
The source text below is rendered literally where necessary; damaged formulas and composite extracts are not silently repaired into a new mathematical claim.
The supplied catalogue marks this record as unrecoverable. Its extract and terminology are documented, but a complete original problem has not been reconstructed.
\paragraph{Statement recorded in the supplied source.}
What about tests for larger dimensions? Are there tradeoffs between the di- mension guarantee, the complexity of the test, and the test’s robustness? Theorem 1.3 (Corollary of MIP* = RE). There exist finite sets X , A and a test T for quantum functions F : X \textbackslash{}to A such that (1) the test T makes only two queries to F , (2) if a quantum function F is finite dimensional, then Pr[T accepts F ] \textbackslash{}leq 21 , (3) there exists a quantum function F such that \href{https://doi.org/10.4171/owr/2021/54}{[S2]}
\paragraph{Source status.}
The supplied catalogue labels this source record partially solved; that metadata was checked against the supplied snapshot on 2026-10-11. The label is preserved for provenance and is not an independent verification of present open status or mathematical correctness. The catalogue literature assessment is dated 2026-08-22.
\subsection{Short English statement}
The fragment asks about dimension, query complexity, and robustness tradeoffs in quantum-function tests.
\subsection{Sources}
\begin{itemize}
\item[S1] Ulam AI and the UnsolvedMath Contributors, UnsolvedMath: A Curated Collection of Open Mathematics Problems, catalogue entry. Catalogue curation and metadata: CC BY 4.0. The source extract is retained; definitions, English summaries and qualifications are editorial additions.
\url{https://huggingface.co/datasets/ulamai/UnsolvedMath}.
\item[S2] Complexity Theory (2022), Oberwolfach Reports, DOI 10.4171/owr/2021/54, pp. 2955--3005
\url{https://doi.org/10.4171/owr/2021/54}.
\item[S3] MIP star equals RE
\url{https://arxiv.org/abs/2001.04383}.
\item Source-specific Oberwolfach Reports licences, including CC BY-SA 4.0 where applicable, are retained; see THIRD\_PARTY\_NOTICES.md.
\end{itemize}
% END PRESERVED DEFINITION
\subsection{Verbatim supplied catalogue question}
What about tests for larger dimensions? Are there tradeoffs between the di- mension guarantee, the complexity of the test, and the test’s robustness? Theorem 1.3 (Corollary of MIP* = RE). There exist finite sets X , A and a test T for quantum functions F : X \textbackslash{}to A such that (1) the test T makes only two queries to F , (2) if a quantum function F is finite dimensional, then Pr[T accepts F ] \textbackslash{}leq 21 , (3) there exists a quantum function F such that

\section{The damaged extract and the established source claim}
The retained catalogue fragment stops before the last displayed
condition of Theorem1.3. Reading the primary report, printed
page2974, recovers that condition as acceptance probability one.
The source presents it as a corollary of $\mathrm{MIP}^*=\mathrm{RE}$:
a two-query test separating its stated infinite-dimensional model
from finite-dimensional strategies. That deep result is existing
source context, not proved or independently validated in this
manuscript. The catalogue fragment must not be turned into a
newly solved dimension tradeoff. We instead attempt a finite
algebraic witness and determine exactly what it can certify.

\section{All moments leave dimension amplification invisible}
In the finite positive-measurement model, replace a family on
$\mathbb C^d$ by $M'_{a\mid x}=M_{a\mid x}\otimes I_k$ on
$\mathbb C^{dk}$ for any integer $k\ge1$. Positivity and
$\sum_a(M'_{a\mid x})^2=I$ are preserved. Every product has
squared Frobenius norm multiplied by $k$, because
$\operatorname{tr}((C\otimes I_k)(C^*\otimes I_k))
=k\operatorname{tr}(CC^*)$. Normalization by dimension also
multiplies by $k$, so all transcript probabilities, including
adaptive ones, agree exactly. No test in this access model can
force an upper bound on the represented Hilbert-space dimension
of every passing realization. A lower bound on the minimum
possible dimension is compatible with this invariance; exact
represented dimension is a different target.

\section{An exact finite relation system with an exponential witness}
Suppose Hermitian matrices $U_1,\ldots,U_{2m}$ on a nonzero
finite-dimensional space satisfy
\[
 U_i^2=I,\qquad U_iU_j=-U_jU_i\quad(i\ne j),\qquad m\ge1.
\]
Then the dimension is divisible by $2^m$, in particular at least
$2^m$. Define $Z_j=iU_{2j-1}U_{2j}$. Anticommutation shows
$Z_j^*=Z_j$ and $Z_j^2=I$. For distinct $j,k$, moving one
pair of $U$'s past the other requires four sign changes, so the
$Z_j$ commute. Their commuting spectral projectors
\[
 P_\varepsilon=\prod_{j=1}^m(I+\varepsilon_jZ_j)/2,
 \qquad\varepsilon\in\{-1,1\}^m,
\]
are orthogonal and sum to $I$: this follows by distributing the
product and using the two complementary projectors for each
involution. At least one has nonzero range.

Conjugation by $U_{2j-1}$ changes the sign of $Z_j$ and preserves
every $Z_k$ for $k\ne j$, using respectively one and two
anticommutation sign changes. Thus this unitary is a bijection
between the joint eigenspaces with sign vectors differing only
at coordinate $j$. By successive flips all $2^m$ joint
eigenspaces have the same positive dimension. Summing their
dimensions proves the assertion.

The bound is attained. With the explicit Pauli matrices
\[
 X=\begin{pmatrix}0&1\\1&0\end{pmatrix},\quad
 Y=\begin{pmatrix}0&-i\\i&0\end{pmatrix},\quad
 Z=\begin{pmatrix}1&0\\0&-1\end{pmatrix},
\]
put on $m$ qubits
\[
 U_{2j-1}=Z^{\otimes(j-1)}\otimes X\otimes I^{\otimes(m-j)},
 \quad
 U_{2j}=Z^{\otimes(j-1)}\otimes Y\otimes I^{\otimes(m-j)}.
\]
Each is Hermitian and squares to $I$. The two at the same
position anticommute because $XY=-YX$. For two at different
positions, the earlier $X$ or $Y$ encounters one $Z$ at its
position, contributing one minus sign; every other tensor factor
commutes. Thus all distinct pairs anticommute. The ambient
dimension is $2^m$, as required. The case $m=0$ is the empty
relation system with arbitrary positive dimension.

\section{The failed passage from relations to a query tradeoff}
The relation system has $2m$ observables and quadratically many
pair relations, but this does not make it a tester with that many
queries. A sequential measurement oracle supplies probabilities
of products of positive measurement operators. It does not grant
coherent access to the signed observables or to arbitrary
polynomial matrix identities. Although an involution can be
written as the difference of its two spectral projectors,
expanding a relation into signed words is not by itself an
accept/reject experiment with perfect completeness and controlled
soundness. No such reduction is supplied here.
Nor does the exact common-eigenspace proof establish stability
under approximate anticommutation: its exact projectors and
exact unitary bijections would require quantitative replacements.
Finally, the finite relation calculation says nothing about the
infinite-dimensional tracial formulation needed for the source's
stated deep corollary.

\section{Audit and status}
The audit checked normalization under tensoring by an identity,
the signs in the $Z_j$ construction, positivity of the common
rank, and the difference between an algebraic certificate and
an accessible sequential test. The recovered final line of the
source theorem is contextual evidence, not a newly supplied proof
of $\mathrm{MIP}^*=\mathrm{RE}$. No computation or theorem prover
was used. This damaged extract remains unresolved as an original
research target; the finite auxiliary witness is complete but
not claimed as original-target progress.

\section{Process and contributions}
Mathematics and model-assisted adversarial self-audit: OpenAI gpt-6-astra, ultra reasoning, as dispatched for this research batch; exact serving revision not exposed. Codex (GPT-6; exact serving revision not exposed) materially assisted with orchestration and format review. No human mathematical authorship, mathematical edits, or independent human proof checking is asserted.
The coordinating Codex (GPT-6; exact serving revision not exposed) also provided a separate model-assisted written-proof review of this manuscript. This is not independent human verification.
Contribution prompt A was applied to the full definition and supplied source row, then prompt B was applied adversarially to the derived results. The item-specific audit above records the scope checks and rejected extrapolations. Local assembly and structural checks do not establish mathematical correctness. The source catalogue status is preserved in TOP\_PROBLEM and is separate from this attempt status.
\paragraph{Integration review.} Codex (GPT-6; exact serving revision not exposed) reviewed the written arguments and expanded layout and source-reference wrappers. Repository metadata and the production web renderer were checked using Python 3.13 and JavaScript with MathJax 3.2.2. This model-assisted review is not independent human proof verification; no PDF compilation or formal proof checking is asserted.

\section{Completion Estimate}
\noindent\textbf{COMPLETION: 0\%}
The proved material is a rederivation of known or auxiliary facts rather than a substantiated new advance on the original target. The source fragment is incomplete and its immediate primary context states an established deep corollary. The exact finite relation witness has no proved sequential testing implementation or robustness guarantee, so no new dimension/query tradeoff is established.
\section{References}
The complete original Sources subsection above is retained, including catalogue attribution and licence notices. The additional list records primary sources actually consulted for this attempt on 11 October 2026; uncited source links reproduced above were not all rechecked.
\begin{itemize}
\item Henry Yuen, Noncommutative Property Testing, OWR54/2021 pp.2971–2974; full primary source passage read in cached text, and page2972 operator definition visually checked from the primary PDF. The original agenda, stated known results and convention mismatch are distinguished. \url{https://ems.press/content/serial-article-files/46932}.
\end{itemize}
\section{Final status}
\textbf{Unresolved in this attempt.} The source fragment is incomplete and its immediate primary context states an established deep corollary. The exact finite relation witness has no proved sequential testing implementation or robustness guarantee, so no new dimension/query tradeoff is established.
\end{document}
